\documentclass[11pt]{article}
\usepackage[margin=1in]{geometry}
\usepackage{amsmath,amssymb,amsthm}
\usepackage{booktabs}
\usepackage{hyperref}

\theoremstyle{plain}
\newtheorem{thm}{Theorem}
\newtheorem{lem}[thm]{Lemma}
\newtheorem{cor}[thm]{Corollary}
\newtheorem{prop}[thm]{Proposition}
\theoremstyle{definition}
\newtheorem{defn}[thm]{Definition}
\theoremstyle{remark}
\newtheorem{rem}[thm]{Remark}
\newtheorem{conj}[thm]{Conjecture}
\theoremstyle{plain}
\newtheorem{gapthm}{Gap}
\newtheorem{retr}{Retraction}

\newcommand{\bI}{\mathbf{I}}
\newcommand{\bM}{\mathbf{M}}
\newcommand{\bK}{\mathbf{K}}
\newcommand{\bL}{\mathbf{L}}
\newcommand{\bu}{\mathbf{u}}
\newcommand{\bw}{\mathbf{w}}
\newcommand{\At}{\widetilde{A}}
\newcommand{\ord}{\operatorname{ord}}
\newcommand{\Res}{\operatorname{Res}}
\newcommand{\Inv}{\operatorname{Inv}}

\title{The long-root collapse multiplier is type-dependent,\\[2mm]
\large and it is the square of the charge-violating boundary entry\\
$\rho^X-1=q\,w_a^{h}\bigl([K^X(w_a)]_0^{\,1}\bigr)^2$, \quad $h=2$ in type $C$, $h=1$ in type $B$;\\
a retraction of the type $B$ claim of \texttt{2026-10-16-c2}, and a reduction of
Gap \texttt{nonvanH} to $q=0$}
\author{Clio}
\date{17 October 2026}

\makeatletter
%% protocol-header.tex -- PROTOCOL §2.3 page-1 block, as a TEMPLATE not an instance.
%%
%% PROTOCOL §2.3: "Every PDF must state, on its first page: author, date,
%% recipient, a title, and the commit hash of the work-in-progress push it
%% corresponds to.  A PDF with no commit hash cannot be checked against the
%% record and should be treated as a draft opinion, not a result."
%%
%% Written 2026-10-16 because the fault had recurred twice and I had twice
%% fixed only the instance:
%%   - 2026-10-14: the Q473 PDF went to Rick with no recipient and no hash.
%%   - 2026-10-14: the Q460 PDF went to Rick saying "Recipient: Robin Langer".
%%     He noticed, and took the commit hash as authority instead.
%% Both times I repaired the one file and wrote "fix the template, not just the
%% instance" in my notes.  This file is that fix.  \input it; do not retype it.
%%
%% USAGE, immediately after \maketitle:
%%   %% protocol-header.tex -- PROTOCOL §2.3 page-1 block, as a TEMPLATE not an instance.
%%
%% PROTOCOL §2.3: "Every PDF must state, on its first page: author, date,
%% recipient, a title, and the commit hash of the work-in-progress push it
%% corresponds to.  A PDF with no commit hash cannot be checked against the
%% record and should be treated as a draft opinion, not a result."
%%
%% Written 2026-10-16 because the fault had recurred twice and I had twice
%% fixed only the instance:
%%   - 2026-10-14: the Q473 PDF went to Rick with no recipient and no hash.
%%   - 2026-10-14: the Q460 PDF went to Rick saying "Recipient: Robin Langer".
%%     He noticed, and took the commit hash as authority instead.
%% Both times I repaired the one file and wrote "fix the template, not just the
%% instance" in my notes.  This file is that fix.  \input it; do not retype it.
%%
%% USAGE, immediately after \maketitle:
%%   %% protocol-header.tex -- PROTOCOL §2.3 page-1 block, as a TEMPLATE not an instance.
%%
%% PROTOCOL §2.3: "Every PDF must state, on its first page: author, date,
%% recipient, a title, and the commit hash of the work-in-progress push it
%% corresponds to.  A PDF with no commit hash cannot be checked against the
%% record and should be treated as a draft opinion, not a result."
%%
%% Written 2026-10-16 because the fault had recurred twice and I had twice
%% fixed only the instance:
%%   - 2026-10-14: the Q473 PDF went to Rick with no recipient and no hash.
%%   - 2026-10-14: the Q460 PDF went to Rick saying "Recipient: Robin Langer".
%%     He noticed, and took the commit hash as authority instead.
%% Both times I repaired the one file and wrote "fix the template, not just the
%% instance" in my notes.  This file is that fix.  \input it; do not retype it.
%%
%% USAGE, immediately after \maketitle:
%%   \input{protocol-header}
%%   \protocolheader{Rick (grandpa-rick)}{clio-vega/proofs}{60e4ce3}
%%
%% The recipient argument is MANDATORY and has no default on purpose: a default
%% is what produced "Recipient: Robin Langer" on a PDF addressed to Rick.

%% \maketitle CLEARS \@date (and \@author, \@title) in the article class, so a
%% macro called AFTER \maketitle that reads \@date prints an EMPTY row.  Caught
%% 2026-10-16 only because I ran pdftotext on page 1 instead of trusting rc=0:
%% the build was clean and the Date field was blank.  Hence: snapshot the date
%% HERE, at \input time in the preamble, where \@date is still live.
\let\protocol@date\@date

\newcommand{\protocolheader}[3]{%
  \begin{center}
  \fbox{\begin{minipage}{0.92\textwidth}
  \small
  \textbf{PROTOCOL §2.3 record.}\\[1mm]
  \begin{tabular}{@{}ll@{}}
  Author     & Clio \\
  Date       & \protocol@date \\
  Recipient  & #1 \\
  Repository & \texttt{#2} \\
  Commit     & \texttt{#3} \\
  \end{tabular}\\[1mm]
  This PDF corresponds to the pushed commit above.  Where this text and the
  commit disagree, \textbf{the commit is authoritative}.
  \end{minipage}}
  \end{center}
  \vspace{2mm}%
}

%%   \protocolheader{Rick (grandpa-rick)}{clio-vega/proofs}{60e4ce3}
%%
%% The recipient argument is MANDATORY and has no default on purpose: a default
%% is what produced "Recipient: Robin Langer" on a PDF addressed to Rick.

%% \maketitle CLEARS \@date (and \@author, \@title) in the article class, so a
%% macro called AFTER \maketitle that reads \@date prints an EMPTY row.  Caught
%% 2026-10-16 only because I ran pdftotext on page 1 instead of trusting rc=0:
%% the build was clean and the Date field was blank.  Hence: snapshot the date
%% HERE, at \input time in the preamble, where \@date is still live.
\let\protocol@date\@date

\newcommand{\protocolheader}[3]{%
  \begin{center}
  \fbox{\begin{minipage}{0.92\textwidth}
  \small
  \textbf{PROTOCOL §2.3 record.}\\[1mm]
  \begin{tabular}{@{}ll@{}}
  Author     & Clio \\
  Date       & \protocol@date \\
  Recipient  & #1 \\
  Repository & \texttt{#2} \\
  Commit     & \texttt{#3} \\
  \end{tabular}\\[1mm]
  This PDF corresponds to the pushed commit above.  Where this text and the
  commit disagree, \textbf{the commit is authoritative}.
  \end{minipage}}
  \end{center}
  \vspace{2mm}%
}

%%   \protocolheader{Rick (grandpa-rick)}{clio-vega/proofs}{60e4ce3}
%%
%% The recipient argument is MANDATORY and has no default on purpose: a default
%% is what produced "Recipient: Robin Langer" on a PDF addressed to Rick.

%% \maketitle CLEARS \@date (and \@author, \@title) in the article class, so a
%% macro called AFTER \maketitle that reads \@date prints an EMPTY row.  Caught
%% 2026-10-16 only because I ran pdftotext on page 1 instead of trusting rc=0:
%% the build was clean and the Date field was blank.  Hence: snapshot the date
%% HERE, at \input time in the preamble, where \@date is still live.
\let\protocol@date\@date

\newcommand{\protocolheader}[3]{%
  \begin{center}
  \fbox{\begin{minipage}{0.92\textwidth}
  \small
  \textbf{PROTOCOL §2.3 record.}\\[1mm]
  \begin{tabular}{@{}ll@{}}
  Author     & Clio \\
  Date       & \protocol@date \\
  Recipient  & #1 \\
  Repository & \texttt{#2} \\
  Commit     & \texttt{#3} \\
  \end{tabular}\\[1mm]
  This PDF corresponds to the pushed commit above.  Where this text and the
  commit disagree, \textbf{the commit is authoritative}.
  \end{minipage}}
  \end{center}
  \vspace{2mm}%
}

\makeatother

\begin{document}
\maketitle
\protocolheader{Robin (\texttt{langer.robin@gmail.com})}{clio-vega/proofs}{0e59afb}

\begin{abstract}
This continues \texttt{proofs/2026-10-16-c2-longroot-multiplier.tex} (henceforth
[LR]) and corrects it.

[LR] Theorem \texttt{collapse} states that on a long-root divisor
$H_{ab}=\{w_aw_b=1\}$ Motegi's normalised weight functions
(arXiv:2610.15344) degenerate by a universal scalar
$\rho=(1-q)^2/\bigl((1-qw_a^2)(1-qw_a^{-2})\bigr)$, and asserts that
``\emph{the same statement and the same $\rho$ hold in type $B$}'',
verified $48/48$.  \textbf{The multiplier claim is false and is retracted here.}
The verification measured $WB/\mathcal{E}^C$: my engine divided the type $B$
partition function by the type \emph{$C$} normalisation factor.  Motegi defines
$\widetilde{WB}:=WB\cdot(\mathcal{E}^B)^{-1}$ with
$\mathcal{E}^B\neq\mathcal{E}^C$, and
\[
  \mathcal{E}^B/\mathcal{E}^C\big|_{\bu=\bw_{\bM}}
  =\prod_{k\notin\bM}\frac{1-qw_k^2}{1-qw_k}
\]
(Lemma~\ref{lem:EBEC}), whose $\bM$-dependence does not cancel in a collapse
ratio.  The correct type $B$ multiplier is
\[
  \rho^B=\frac{(1-q)^2}{(1-qw_a)(1-qw_b)},
  \qquad\text{against}\qquad
  \rho^C=\frac{(1-q)^2}{(1-qw_a^2)(1-qw_b^2)}
\]
(Theorem~\ref{thm:rhoB}): the \emph{long}-root factors $1-qe^{+2t_c}$ of type
$C$ are replaced by the \emph{short}-root factors $1-qe^{+t_c}$.  Each candidate
holds on its own type and fails on the other, over a domain on which they
disagree everywhere.

The correction is not a loss: it supplies the discriminating coordinate that
makes the boundary reading a statement rather than a resemblance.  Writing
$[K^X(u)]_0^{\,1}$ for the charge-violating entry of Motegi's triangular
$K$-matrix --- $1-u^{-2}$ in type $C$, $1-u^{-1}$ in type $B$ ---
Theorem~\ref{thm:dict} is the identity
\[
  \rho^X-1=\frac{q\,w_a^{h}\bigl([K^X(w_a)]_0^{\,1}\bigr)^2}
                {(1-qw_a^{h})(1-qw_b^{h})},
  \qquad h=2\ (X=C),\quad h=1\ (X=B),
\]
with $h$ the length of the boundary root.  The naive form of this reading uses
the \emph{unnormalised} entry, which is $\beta,\gamma$-free hence identical in
both types, and therefore predicts a type-independent $\rho$: that form is
refuted by Theorem~\ref{thm:rhoB}.  The two readings agree at $w_a=+1$ and
disagree at $w_a=-1$, and the measurement lands on the disagreement set.  The
slot is also pinned, and from a source theorem rather than from my own algebra:
Motegi's divisibility theorems make
$\prod_{m\in\bM}\bigl(1-w_m^{-h}\bigr)$ a divisor of
$\widetilde{WX}_{\bM,\bI}$, indexed by the specialisation set $\bM$ --- which is
exactly the index set of the $K$-matrices, since
$\At_{\bM,\bI}$ is the specialisation $\bu=\bw_{\bM}$ and $K$ acts on the
$\bu$-slot.  A long-root collapse deletes $\{a,b\}$ from $\bM$, hence deletes
exactly the two $K$-matrices $K(w_a)$, $K(w_b)$.

On the primary target, Gap \texttt{nonvanH} ($\ord_{H_{ab}}\At_{\bM,\bI}=0$ when
$\bM\ge\bI$ and $\{a,b\}\cap\bM=\emptyset$), I report a reduction and two failed
attempts with their breakdown points.  The reduction
(Proposition~\ref{prop:q0}) is that it suffices to prove the statement at
$q=0$, because $\At_{\bM,\bI}$ is regular at $q=0$ and $\mathcal{E}^C|_{q=0}=1$;
and that at $q=0$ the entire dependence of $\At_{\bM,\bI}$ on the variables
outside $\bw_{\bM}$ is the single product $\prod_{j\le r}\prod_{k<i_j}(u_j-w_k)$,
so that $w_k$ is absent whenever $k\ge\max\bI$.  This closes the gap on
the stratum $\max\bI\le b$ --- $100\%$, $91.7\%$, $81.9\%$, $72.8\%$, $64.8\%$,
$58.0\%$, $52.3\%$ of the domain at $N=2,\dots,8$, \emph{decreasing} --- modulo
nonvanishing of $\At_{\bM,\bI}$ at $q=0$, which is a factorial $GQ$-function
statement and is \emph{stronger} than the converse of triangularity that
[Q488] already carries as measured.  I say so rather than calling it a closure.
The two failed attempts are recorded in \S\ref{sec:failed}: a bidegree-corner
argument, \emph{refuted by measurement} (the corner is empty in $12/12$ cases at
$N=3$), and its $q=0$ refinement, which reduces to a symmetrisation sum twisted
by a non-symmetric monomial that I cannot evaluate.  Gap \texttt{nonvanH}
remains open.

Everything [LR] states for type $C$ is untouched: $\mathcal{E}^C$ is the correct
normaliser for type $C$ and is checked against Motegi's independent diagonal
root-product formula.
\end{abstract}

\tableofcontents

\section{What this corrects, and how the error was found}
\label{sec:intro}

The session brief armed a gate at position zero: a dictionary, to be written by
hand, deciding whether the variable $w_a$ appearing in
\[
  \rho-1=\frac{q\,(w_a-w_a^{-1})^2}{(1-qw_a^2)(1-qw_a^{-2})}
  \tag{[LR] Prop.\ \texttt{rho}}
\]
occupies the same slot as the argument $u$ of Motegi's boundary $K$-matrix.  The
brief's own pre-registered worry was the right one --- ``same symbol, possibly
different slot'' --- and the gate paid for itself three times over, in ways the
brief did not anticipate.

\paragraph{First: the brief quoted the wrong paper.}
The brief records, as verified verbatim at wake,
${}_a\langle1|K_a(u)|0\rangle_a=u-u^{-1}$ from
\texttt{arxiv-1703.07924-source.TEX}, eq.\ \texttt{kmatrix}, L300--321.  That
citation is accurate, and it is about a different paper.  The theorem under
discussion is about arXiv:2610.15344, whose triangular $K$-matrix carries the
\emph{opposite} triangularity (eq.\ \texttt{eq:Ktri-components}, L332--345):
\begin{equation}
[K(u)]_i^{\,j}=
\begin{cases}
\beta-\gamma u^{-1} & (i,j)=(0,0),\\
1-u^{-2}            & (i,j)=(0,1),\\
0                   & (i,j)=(1,0),\\
\beta u^{-2}-\gamma u^{-1} & (i,j)=(1,1).
\end{cases}
\label{eq:Ktri}
\end{equation}
In 1703.07924 the charge-violating entry is the $(1,0)$ entry and equals
$u-u^{-1}$; here the $(1,0)$ entry is \emph{zero} and the charge-violating entry
is the $(0,1)$ entry, $1-u^{-2}$.  The two are proportional,
$u-u^{-1}=u\,(1-u^{-2})$, so the character-level match that excited the brief
survives --- but only one of the two is a quantity of the paper being proved
about, and the discrepancy by a factor of $u$ is exactly the monomial $w_a^h$
that appears in Theorem~\ref{thm:dict}.  Two Motegi papers, the same phrase
``triangular $K$-matrix'', transposed conventions.

\paragraph{Second: the two specialisations of \eqref{eq:Ktri}.}
Motegi specialises \eqref{eq:Ktri} twice
(eqs.\ \texttt{eq:KC-components}, \texttt{eq:KB-components}, \texttt{triandc},
\texttt{triandb}, L352--375):
\begin{equation}
K^C(u)=K(u)\big|_{\gamma=0,\beta=1},
\qquad
K^B(u)=(1+u^{-1})^{-1}K(u)\big|_{\gamma=-1,\beta=1},
\label{eq:KCKB}
\end{equation}
so that the charge-violating entries are
\begin{equation}
[K^C(u)]_0^{\,1}=1-u^{-2},
\qquad
[K^B(u)]_0^{\,1}=1-u^{-1}.
\label{eq:Kentries}
\end{equation}
Note that in \eqref{eq:Ktri} the $(0,1)$ entry does not involve $\beta$ or
$\gamma$ at all.  So the \emph{unnormalised} charge-violating entry is literally
the same function in both types, and only the type $B$ scalar $(1+u^{-1})^{-1}$
of \eqref{eq:KCKB} distinguishes them.  This is the fork on which the whole
reading turns, and \S\ref{sec:dict} resolves it.

\paragraph{Third: the engine was normalising type $B$ by $\mathcal{E}^C$.}
Checking the slot required reading Motegi's definitions of the two normalisation
factors, and they are not the same function.  That is \S\ref{sec:EB}.  It
invalidates every type $B$ number in [LR] and in [Q484], and it changes a
theorem.

\paragraph{Notation.}
As in [LR].  $\mathbb{F}=\mathbb{C}(q)(w_1,\dots,w_N)$; for
$\bM,\bI\subseteq[N]$ with $|\bI|=r$ and $|\bM|=s\ge r$,
\begin{equation}
\At^X_{\bM,\bI}:=\widetilde{WX}_{N,s,r}(\bw_{\bM}\,|\,\bw\,|\,\bI),
\qquad X\in\{B,C\},
\label{eq:Atdef}
\end{equation}
and $\At_{\bM,\bI}:=\At^C_{\bM,\bI}$.  Throughout $a<b$ are fixed,
$H:=H_{ab}=\{w_aw_b=1\}$, $t:=w_aw_b-1$, and $f\mapsto f|_H$ is restriction
(i.e.\ the substitution $w_b\mapsto w_a^{-1}$).  Motegi's order $\bM\ge\bI$ is
as restated in [LR] Lemma \texttt{order}.  All line numbers refer to the e-print
source \texttt{typebcweightfunctions\_arxivversionfinal-3.tex} of
arXiv:2610.15344 unless the file is named otherwise.

\section{The two normalisation factors}
\label{sec:EB}

Motegi's unnormalised weight functions are two specialisations of one partition
function $F_{N,n,r}(\bu|\bw|\bI)$ (eqs.\ \texttt{WCfunction},
\texttt{WBfunction}, L489--505):
\begin{equation}
WC=F\big|_{\gamma=0,\beta=1},
\qquad
WB=\Bigl(\prod_{k=1}^n(1+u_k^{-1})^{-1}\Bigr)F\big|_{\gamma=-1,\beta=1},
\label{eq:WCWB}
\end{equation}
and the normalised ones are
$\widetilde{WX}:=WX\cdot(\mathcal{E}^X)^{-1}$
(eqs.\ \texttt{typeCnormalizeddefinition} L1051,
\texttt{typeBnormalizeddefinition} L1358), with
\begin{align}
\mathcal{E}^C_{N,n,r}&=\frac{1}{(1-q)^{n-r}}\,
\frac{\prod_{j,k=1}^n(1-qu_j^{-1}u_k)\ \prod_{j=1}^n\prod_{k=1}^N(1-qu_jw_k)}
     {\prod_{1\le j<k\le n}(1-qu_ju_k)\ \prod_{1\le j\le k\le N}(1-qw_jw_k)},
\label{eq:EC}\\[1mm]
\mathcal{E}^B_{N,n,r}&=\frac{1}{(1-q)^{n-r}}\,
\frac{\prod_{j,k=1}^n(1-qu_j^{-1}u_k)\ \prod_{j=1}^n\prod_{k=1}^N(1-qu_jw_k)\
      \prod_{j=1}^n(1-qu_j)}
     {\prod_{1\le j<k\le n}(1-qu_ju_k)\ \prod_{j=1}^n(1-qu_j^2)\
      \prod_{1\le j<k\le N}(1-qw_jw_k)\ \prod_{k=1}^N(1-qw_k)}
\label{eq:EB}
\end{align}
(eqs.\ \texttt{typeCnormalizationfactor} L1057--1062,
\texttt{typeBnormalizationfactor} L1361--1367).

The difference is not cosmetic and it is not a scalar.  Read the two
denominators: $\mathcal{E}^C$ carries the $\bw$-product over $j\le k$, i.e.\
\emph{including} the diagonal $\prod_k(1-qw_k^2)$, and no $\bu$-diagonal;
$\mathcal{E}^B$ carries the $\bw$-product over $j<k$ \emph{strictly}, has the
diagonal moved to the $\bu$-side as $\prod_j(1-qu_j^2)$, and in its place
carries the linear product $\prod_k(1-qw_k)$.  Under the specialisation
$\bu=\bw_{\bM}$ the $\bu$-diagonal becomes a $\bw$-diagonal over $\bM$, and the
two conventions then differ exactly on the complement.

\begin{lem}[the ratio of the normalisers]\label{lem:EBEC}
For $\bM\subseteq[N]$ with $|\bM|=n$, and any $r\le n$,
\begin{equation}
\frac{\mathcal{E}^B_{N,n,r}}{\mathcal{E}^C_{N,n,r}}\bigg|_{\bu=\bw_{\bM}}
=\prod_{k\in[N]\setminus\bM}\frac{1-qw_k^2}{1-qw_k}.
\label{eq:EBEC}
\end{equation}
Equivalently, with $G_{\bM}$ denoting the right-hand side,
$\At^{B}_{\bM,\bI}=\bigl(WB/\mathcal{E}^C\bigr)\big|_{\bu=\bw_{\bM}}\cdot
G_{\bM}^{-1}$.
\end{lem}

\begin{proof}
Every factor of \eqref{eq:EC} occurs in \eqref{eq:EB}, so the quotient is
\[
\frac{\mathcal{E}^B}{\mathcal{E}^C}
=\frac{\prod_{j=1}^n(1-qu_j)}{\prod_{j=1}^n(1-qu_j^2)}\cdot
 \frac{\prod_{1\le j\le k\le N}(1-qw_jw_k)}
      {\prod_{1\le j<k\le N}(1-qw_jw_k)\ \prod_{k=1}^N(1-qw_k)}
=\prod_{j=1}^n\frac{1-qu_j}{1-qu_j^2}\cdot
 \prod_{k=1}^N\frac{1-qw_k^2}{1-qw_k},
\]
the second equality because the two $\bw$-products over pairs cancel except on
the diagonal.  Substituting $\bu=\bw_{\bM}$ makes the first product
$\prod_{m\in\bM}\frac{1-qw_m}{1-qw_m^2}$, which cancels the $\bM$-part of the
second, leaving \eqref{eq:EBEC}.  Note that no factor simplifies further:
$1-qu^2$ is not $(1-qu)(1+qu)$, the $q$ being to the first power in each.
\end{proof}

\begin{rem}[the swap is the long/short root swap]\label{rem:rootswap}
Lemma~\ref{lem:EBEC} is the statement that type $C$'s ascent factor at the index
$k$ is $1-qw_k^2=1-q\,e^{+2t_k}$ and type $B$'s is $1-qw_k=1-q\,e^{+t_k}$:
the \emph{long} root $2t_k$ of $C_N$ against the \emph{short} root $t_k$ of
$B_N$.  This is corroborated, independently of \eqref{eq:EC}--\eqref{eq:EB}, by
Motegi's two divisibility theorems (eqs.\ \texttt{typeCfullDivisibility}
L1086--1089 and \texttt{typeBfullDivisibility} L1394--1402), which state that
for $|\bM|=|\bI|=r$ the function $\widetilde{WX}_{\bM,\bI}$ is divisible by
\begin{center}
\begin{tabular}{@{}lll@{}}
\toprule
 & $\bM$-side (inversion) & $\bM^{\mathrm{c}}$-side (ascent) \\
\midrule
$X=C$ & $\prod_{1\le j\le k\le r}(1-w_{m_j}^{-1}w_{m_k}^{-1})$
      & $\prod_{j\le k\in\bM^{\mathrm{c}}}(1-qw_jw_k)$ \\
$X=B$ & $\prod_{j}(1-w_{m_j}^{-1})\prod_{j<k}(1-w_{m_j}^{-1}w_{m_k}^{-1})$
      & $\prod_{j\in\bM^{\mathrm{c}}}(1-qw_j)\prod_{j<k\in\bM^{\mathrm{c}}}(1-qw_jw_k)$ \\
\bottomrule
\end{tabular}
\end{center}
(both also carry $\prod_{j\in\bM^{\mathrm{c}}}\prod_{l:m_l<j}(1-qw_{m_l}^{-1}w_j)$).
The diagonals are $1-w_m^{-2}$ versus $1-w_m^{-1}$ on the $\bM$-side and
$1-qw_j^2$ versus $1-qw_j$ on the complement: the same swap, on both sides, in a
theorem that is not a normalisation factor.  This is what \S\ref{sec:dict}
consumes, and it is how the implementation of \eqref{eq:EB} is checked
(\S\ref{sec:verif}).
\end{rem}

\section{The type \texorpdfstring{$B$}{B} collapse multiplier}
\label{sec:rhoB}

\begin{retr}\label{retr:typeB}
[LR] Theorem \texttt{collapse} ends with: ``\emph{The same statement and the
same $\rho$ hold in type $B$}'', and [LR] \S\texttt{verif} reports
$48/48$ at $N=3$ for type $B$ with $12$ non-degenerate.  The \emph{statement}
survives; the words ``the same $\rho$'' are withdrawn.  The computation behind
the $48/48$ divided $WB$ by $\mathcal{E}^C$ --- the engine
\texttt{q484\_engine.Atilde} carries an in-code comment asserting that
``$\mathcal{E}^C$ is the common normalizer'', which is false --- so it is a
correct measurement of $WB/\mathcal{E}^C$ and not of
$\widetilde{WB}=WB/\mathcal{E}^B$.  By the same token the type $B$ arm of
[Q484] Theorem A ($48/48$) measures the same wrong object and is withdrawn
pending re-measurement.  Nothing stated for type $C$ is affected:
$\mathcal{E}^C$ is the correct normaliser there.
\end{retr}

\begin{thm}[the type $B$ multiplier]\label{thm:rhoB}
Let $a<b$, $\{a,b\}\subseteq\bM$, $\bM'=\bM\setminus\{a,b\}$.  Then for every
$\bI\subseteq[N]$,
\begin{equation}
\At^{B}_{\bM,\bI}\equiv\rho^B\cdot\At^{B}_{\bM',\bI}\pmod{w_aw_b-1},
\qquad
\rho^B=\frac{(1-q)^2}{(1-qw_a)(1-qw_b)},
\label{eq:rhoB}
\end{equation}
and $\rho^B\ne\rho^C$, where
$\rho^C=(1-q)^2/\bigl((1-qw_a^2)(1-qw_b^2)\bigr)$ is the multiplier of [LR]
Theorem \texttt{collapse} (there written with $w_b^2$ replaced by its value
$w_a^{-2}$ on $H$).  Uniformly in the type,
\begin{equation}
\rho^X=\frac{(1-q)^2}{(1-qw_a^{h})(1-qw_b^{h})},
\qquad h=2\ (X=C),\quad h=1\ (X=B),
\label{eq:rhoX}
\end{equation}
with $h$ the length of the boundary root: $2t_c$ in type $C$, $t_c$ in type $B$.
\end{thm}

\begin{proof}
Write $\Lambda_{\bM,\bI}:=\bigl(WB/\mathcal{E}^C\bigr)|_{\bu=\bw_{\bM}}$ for the
object [LR] measured, so that
$\At^{B}_{\bM,\bI}=\Lambda_{\bM,\bI}\,G_{\bM}^{-1}$ by
Lemma~\ref{lem:EBEC}, with
$G_{\bM}=\prod_{k\notin\bM}\frac{1-qw_k^2}{1-qw_k}$.  Then
\begin{equation}
\frac{\At^{B}_{\bM,\bI}}{\At^{B}_{\bM',\bI}}
=\frac{\Lambda_{\bM,\bI}}{\Lambda_{\bM',\bI}}\cdot\frac{G_{\bM'}}{G_{\bM}},
\qquad
\frac{G_{\bM'}}{G_{\bM}}
=\prod_{k\in\bM'^{\mathrm{c}}\setminus\bM^{\mathrm{c}}}\frac{1-qw_k^2}{1-qw_k}
=\frac{(1-qw_a^2)(1-qw_b^2)}{(1-qw_a)(1-qw_b)},
\label{eq:Gratio}
\end{equation}
since $\bM'^{\mathrm{c}}=\bM^{\mathrm{c}}\sqcup\{a,b\}$.  [LR]'s type $B$
measurement, reproduced here (\S\ref{sec:verif}), is exactly
$\Lambda_{\bM,\bI}|_H=\rho^C\,\Lambda_{\bM',\bI}|_H$, and the quotient
\eqref{eq:Gratio} is a unit at $H$ with no pole or zero there (each factor is
$1-q\cdot(\text{monomial})$, nonzero because $q$ is an indeterminate).
Substituting,
\[
\At^{B}_{\bM,\bI}\big|_H
=\rho^C\cdot\frac{(1-qw_a^2)(1-qw_b^2)}{(1-qw_a)(1-qw_b)}\cdot
 \At^{B}_{\bM',\bI}\big|_H
=\frac{(1-q)^2}{(1-qw_a)(1-qw_b)}\,\At^{B}_{\bM',\bI}\big|_H .
\]
That $\rho^B\neq\rho^C$ in $\mathbb{F}$ is \eqref{eq:Gratio} together with
$(1-qw_a^2)(1-qw_b^2)\neq(1-qw_a)(1-qw_b)$.
\end{proof}

\begin{rem}[where the proof of {[LR]} Theorem \texttt{collapse} goes through, and
where it does not]\label{rem:reproof}
The proof in [LR] is a factor-by-factor comparison of the branching formula at
$\bM$ and at $\bM'$, and its final step is: ``\emph{what survives is the one
factor of $\Pi_1$ with both legs in $\{a,b\}$, against the leading
$(1-q)^2$}'', where $\Pi_1=\prod_{j\le k\in\bM^{\mathrm{c}}}(1-qw_jw_k)$ is the
image in the numerator of $\At$ of the $\bw$-product in the denominator of
$\mathcal{E}^C$.  In type $B$ that product is
$\prod_{j<k\in\bM^{\mathrm{c}}}(1-qw_jw_k)\cdot\prod_{j\in\bM^{\mathrm{c}}}(1-qw_j)$
--- strictly off-diagonal, plus a linear product (Remark~\ref{rem:rootswap}).
So there \emph{is} no factor with both legs in $\{a,b\}$ beyond
$(1-qw_aw_b)$, which restricts to $1-q$ and is cancelled by $\Pi_4$'s
$(1-qw_a^{-1}w_b^{-1})|_H=1-q$ exactly as in [LR]; and what survives instead is
the pair of linear factors $(1-qw_a)(1-qw_b)$.  The proof above routes around
this by going through $\Lambda$ rather than re-deriving the branching formula in
type $B$, which is why it does not need the type $B$ cross-sector branching
formula --- an input [LR] asserted (``the cross-sector branching formula has the
same shape'') without checking.  The factor-by-factor reading and the
$\Lambda$-route give the same answer, which is one of the two independent
confirmations recorded in \S\ref{sec:verif}.

The sentence in [LR] that \emph{was} right, and which now reads as the warning it
should have been, is: ``the type $B$ specialisation differs from type $C$ by
$\beta=1,\gamma=-1$ and the prefactor $\prod_k(1+u_k^{-1})^{-1}$, neither of
which contributes a factor vanishing on a long-root divisor.''  Both halves of
that are true.  Neither bears on the \emph{value} of the multiplier, and the
normalisation factor --- the thing that actually differs --- is not mentioned.
A correct observation about vanishing was allowed to stand in for an unmade
observation about value.
\end{rem}

\section{The dictionary: \texorpdfstring{$\rho^X-1$}{rho-1} is a boundary entry, squared}
\label{sec:dict}

\begin{prop}\label{prop:rhoXminus1}
With $h$ as in \eqref{eq:rhoX}, on $H$ (where $w_b=w_a^{-1}$),
\begin{equation}
\rho^X-1
=\frac{q\,\bigl(w_a^{h/2}-w_a^{-h/2}\bigr)^2}{(1-qw_a^{h})(1-qw_a^{-h})},
\label{eq:rhoXminus1}
\end{equation}
so that $\ord_H\rho^X=\ord_H(\rho^X-1)=0$ and $\rho^X\ne1$.  Explicitly, the
numerator is $q(w_a-w_a^{-1})^2$ in type $C$ and $q(w_a+w_a^{-1}-2)$ in type
$B$.
\end{prop}
\begin{proof}
$(1-q)^2-(1-qx)(1-qx^{-1})=(1-2q+q^2)-(1-q(x+x^{-1})+q^2)
=q(x+x^{-1}-2)=q(x^{1/2}-x^{-1/2})^2$ for any $x$; apply it with
$x=w_a^{h}$, using $w_b^{h}=w_a^{-h}$ on $H$.  Neither numerator nor denominator
involves $w_b$, so neither is divisible by $t$.
\end{proof}

\begin{thm}[the boundary dictionary]\label{thm:dict}
For $X\in\{B,C\}$, with $[K^X(u)]_0^{\,1}$ the charge-violating entry
\eqref{eq:Kentries} of Motegi's $K$-matrix and $h=2$ resp.\ $1$,
\begin{equation}
[K^X(u)]_0^{\,1}=u^{-h/2}\bigl(u^{h/2}-u^{-h/2}\bigr),
\qquad\text{hence}\qquad
\rho^X-1=\frac{q\,w_a^{h}\,\bigl([K^X(w_a)]_0^{\,1}\bigr)^2}
              {(1-qw_a^{h})(1-qw_b^{h})}\bigg|_H .
\label{eq:dict}
\end{equation}
In particular the zero locus in $w_a$ of the obstruction $\rho^X-1$ equals the
zero locus of the charge-violating boundary entry, in each type separately:
$\{w_a=\pm1\}$ in type $C$ and $\{w_a=+1\}$ in type $B$.
\end{thm}
\begin{proof}
$u^{-1}(u-u^{-1})=1-u^{-2}=[K^C(u)]_0^{\,1}$ and
$u^{-1/2}(u^{1/2}-u^{-1/2})=1-u^{-1}=[K^B(u)]_0^{\,1}$, which is the first
identity; substituting it into \eqref{eq:rhoXminus1} gives the second.  Both
were also checked symbolically (\S\ref{sec:verif}).  For the loci:
$1-u^{-2}$ vanishes at $u=\pm1$ and $1-u^{-1}$ only at $u=+1$; and
$q(w_a-w_a^{-1})^2$ vanishes at $w_a=\pm1$ while
$q(w_a+w_a^{-1}-2)=q\,w_a^{-1}(w_a-1)^2$ takes the value $-4q\neq0$ at
$w_a=-1$.
\end{proof}

\subsection{Why this is a statement and not a resemblance}

The identity \eqref{eq:dict} has to be defended against the objection that a
two-variable polynomial of the shape $q(x^{1/2}-x^{-1/2})^2$ can be factored to
look like anything.  Three coordinates, of which the brief had checked one.

\paragraph{Coordinate 1: the polynomial.}
Checked by the brief, and it was checked against the wrong paper
(\S\ref{sec:intro}).  With the right entry the match is better, not worse: the
discrepancy of a factor $u$ between $u-u^{-1}$ and $1-u^{-2}$ is the monomial
$w_a^{h}$ of \eqref{eq:dict}, and the brief's reading of
$(w_a-w_a^{-1})^2$ as a \emph{square} is correct --- the square of one entry at
one argument, not a product of two entries at two arguments, and not the
kernel-argument doubling of the Pfaffian literature that the brief feared.

\paragraph{Coordinate 2: the slot, the index set, and the count.}
Motegi's $K$-matrices act on the $\bu$-slot, one per reflected row, with
arguments $u_1,\dots,u_n$ (eqs.\ \eqref{eq:WCWB}, \texttt{eq:reflection} L382,
and ``$K_a(u)$ denotes the operator acting non-trivially on the $a$-th
component'' L388).  By \eqref{eq:Atdef}, $\At_{\bM,\bI}$ \emph{is} the
specialisation $\bu=\bw_{\bM}$.  Hence:
\begin{center}
\emph{the $K$-matrices of $\At_{\bM,\bI}$ are exactly $\{K(w_m):m\in\bM\}$,
indexed by the specialisation set $\bM$} ---
\end{center}
not by the label $\bI$, and not by the complement.  A long-root collapse deletes
$\{a,b\}$ from $\bM$ (the hypothesis of Theorem~\ref{thm:rhoB} is
$\{a,b\}\subseteq\bM$), hence deletes exactly two $K$-matrices, $K(w_a)$ and
$K(w_b)$, and $\rho^X$ is by construction the ratio of the two levels.  Two rows
deleted, two boundary entries, at the two arguments $w_a$, $w_b$ --- and on $H$
these are mutually inverse, $w_b=w_a^{-1}$, which is why one squared entry at
one argument suffices.

This is not read off my own algebra.  Remark~\ref{rem:rootswap} records that
Motegi's divisibility theorems make
$\prod_{m\in\bM}\bigl(1-w_m^{-h}\bigr)
=\prod_{m\in\bM}[K^X(w_m)]_0^{\,1}$
an explicit divisor of $\widetilde{WX}_{\bM,\bI}$, with the product running over
$\bM$ and with the type-correct exponent $h$ in each case.  The charge-violating
boundary entry, indexed by the $K$-argument set, is a divisor of the weight
function in the source.

\paragraph{Coordinate 3: the disagreement set.}
Here is the discrimination, and it is why Theorem~\ref{thm:rhoB} was worth the
retraction.  There are two candidate readings:
\begin{itemize}
\item[(U)] the \emph{unnormalised} entry $[K(u)]_0^{\,1}=1-u^{-2}$ of
  \eqref{eq:Ktri}, which is $\beta,\gamma$-free and therefore the same function
  in both types.  Reading (U) predicts a \emph{type-independent} $\rho$.
\item[(N)] the \emph{normalised} entries $[K^X(u)]_0^{\,1}$ of
  \eqref{eq:Kentries}, which differ between the types by precisely the type $B$
  scalar $(1+u^{-1})^{-1}$.  Reading (N) predicts a type-dependent $\rho$,
  tracking $h$.
\end{itemize}
The two readings \emph{agree} at $w_a=+1$ and \emph{disagree} at $w_a=-1$,
which is the only point where the type $B$ scalar $(1+u^{-1})^{-1}$ is singular.
The measurement lands on the disagreement set:
\begin{center}
\begin{tabular}{@{}lcccc@{}}
\toprule
 & $[K^C]_0^{\,1}$ & $\rho^C-1$ & $[K^B]_0^{\,1}$ & $\rho^B-1$ \\
\midrule
$w_a=+1$ & $0$ & $0$ & $0$ & $0$ \\
$w_a=-1$ & $0$ & $0$ & $2$ & $-4q\ \neq0$ \\
\bottomrule
\end{tabular}
\end{center}
Reading (U) is refuted and reading (N) is confirmed, on a locus where they make
opposite predictions.  Had [LR]'s type $B$ claim been true, (U) would have been
the surviving reading and the identity would have carried no information about
the boundary at all --- a $\beta,\gamma$-free entry cannot distinguish anything.
It is the \emph{error} that supplied the discriminating coordinate.

\begin{rem}[the point the brief pre-registered]\label{rem:task2}
The brief's Task 2 pre-registered a falsifier at $w_a=w_b=\pm1$, where
$w_aw_b=1$ and $w_a^{h}=1$ both hold, and named the right arm for it:
``\emph{the type $B$ $WB$ prefactor $\prod_k(1+u_k^{-1})^{-1}$ has a pole at
$u=-1$ and is not $u\mapsto-u$ invariant, so $w_a=-1$ is exactly where $B$ can
differ from $C$ --- same point, reached from the other arm.  Run both signs
separately.}''  Run both signs separately: that is the table above, and the
instruction was exactly right.  The brief expected the falsifier to kill the
reading; it killed one of the reading's two forms and promoted the other.
\end{rem}

\begin{rem}[what is still not established]\label{rem:coord3}
The \emph{occurrence} of $([K^X]_0^{\,1})^2$ in $\rho^X-1$ is not shown to be
\emph{causal}.  In the derivation the numerator of $\rho^X$ is $(1-q)^2$, which
comes from the bulk factor $(1-q)^{|\bL|}$ of the branching formula and contains
no $K$-matrix; the boundary entry appears only after subtracting $1$ and
restricting to $H$.  So the status of the reading is: polynomial matched, slot
and index set and count and type-dependence matched, mechanism unproved.  It is
mathematics with a named mechanism gap, not a theorem about the lattice model.
I record this because the identity is pretty and the temptation to call it an
explanation is the exact failure the gate exists to catch.
\end{rem}

\section{Gap \texttt{nonvanH}: a reduction, a stratum, and two failures}
\label{sec:nonvanH}

The primary target of the session was [LR]'s remaining named hypothesis.

\begin{gapthm}[\texttt{nonvanH}]\label{gap:nonvanH}
$\ord_{H_{ab}}\At_{\bM,\bI}=0$ whenever $\bM\ge\bI$ and
$\{a,b\}\cap\bM=\emptyset$.
\end{gapthm}

Equivalently, since $\At_{\bM,\bI}$ is a Laurent polynomial (fact (R) of [LR]):
$t=w_aw_b-1$ does not divide $\At_{\bM,\bI}$.  It remains open.  What follows is
a stratum check, a reduction, a closure on part of the domain, and two attempts
with their breakdown points.

\subsection{Stratum check: nothing is borrowable from the collapse}
\label{sec:strat}

[LR] Theorem \texttt{collapse} has hypothesis $\{a,b\}\subseteq\bM$;
Gap~\ref{gap:nonvanH} has hypothesis $\{a,b\}\cap\bM=\emptyset$.  These strata
are \emph{complementary}, so the collapse theorem says nothing whatever on the
domain of the gap, and no part of the proof below may borrow from it.  The
degenerate reading confirms that they must be disjoint: if \eqref{eq:rhoB} held
with $\{a,b\}\cap\bM=\emptyset$ then $\bM'=\bM$ and it would read
$(\rho-1)\At_{\bM,\bI}\equiv0\pmod t$, whence $\At_{\bM,\bI}|_H=0$ since
$\rho-1$ is a unit at $H$ (Proposition~\ref{prop:rhoXminus1}) --- the exact
negation of Gap~\ref{gap:nonvanH}.  The two statements are about disjoint strata
and both can be true; on the overlap they could not be.

For completeness: the third stratum $|\{a,b\}\cap\bM|=1$ is not needed.  The
hypothesis \eqref{eq:hyp} of [LR] Theorem \texttt{created} is applied only at
$\bM=\nu'=\nu\setminus\{a,b\}$, which is disjoint from $\{a,b\}$ by
construction; and the stratum $\{a,b\}\subseteq\bM$ reduces to the disjoint one
by the collapse theorem, since $\rho$ is a unit at $H$.

\subsection{Reduction to \texorpdfstring{$q=0$}{q=0}}
\label{sec:q0}

\begin{lem}\label{lem:q0reg}
$\At_{\bM,\bI}$ is regular at $q=0$, and $\mathcal{E}^C_{N,n,r}|_{q=0}=1$;
consequently $\At_{\bM,\bI}|_{q=0}
=F_{N,n,r}(\bw_{\bM}|\bw|\bI)\big|_{q=0,\ \gamma=0,\beta=1}$.
\end{lem}
\begin{proof}
Every factor of \eqref{eq:EC} other than the prefactor has the form
$1-q\cdot(\text{monomial})$, hence equals $1$ at $q=0$; and the prefactor
$(1-q)^{-(n-r)}$ equals $1$ there too.  So
$\mathcal{E}^C|_{q=0}=1$, and in particular $\mathcal{E}^C$ is a unit of the
local ring at $q=0$, so $(\mathcal{E}^C)^{-1}$ is regular there.  $F$ is
polynomial in $q$ by Motegi's Theorem 3.3 (L491--504; the only denominators are
the $u_j-u_k$, free of $q$).  Hence $\At=F\cdot(\mathcal{E}^C)^{-1}$ is regular
at $q=0$ with the stated value.
\end{proof}

\begin{prop}[the reduction]\label{prop:q0}
If $\At_{\bM,\bI}\big|_{q=0}\big|_H\ne0$ then $\At_{\bM,\bI}\big|_H\ne0$.  Hence
it suffices to prove Gap~\ref{gap:nonvanH} at $q=0$.
\end{prop}
\begin{proof}
By Lemma~\ref{lem:q0reg} the two operations $f\mapsto f|_{q=0}$ and
$f\mapsto f|_H$ are both defined on $\At_{\bM,\bI}$ and commute (they
substitute into disjoint sets of variables, $q$ and $w_b$).  So
$\At|_H=0$ forces $\At|_{q=0}|_H=0$.
\end{proof}

The reduction is worth having because $q=0$ collapses the formula almost
completely.

\begin{lem}[the $q=0$ structure]\label{lem:q0struct}
At $q=0$ the entire dependence of $\At_{\bM,\bI}$ on the variables
$\{w_k:k\notin\bM\}$ is through the single product
$\prod_{j=1}^{r}\prod_{k<i_j}(u_j-w_k)$, with $\bu=\bw_{\bM}$.  Explicitly,
\begin{equation}
\At_{\bM,\bI}\big|_{q=0}
=\frac{\prod_{k=1}^{n}u_k^{-(N+n)}}{(n-r)!}
 \sum_{\sigma\in S_n}\sigma\bigl[\mathcal{B}^0\bigr],
\label{eq:q0expand}
\end{equation}
\begin{equation}
\mathcal{B}^0=\prod_{j=1}^{r}\Bigl[(u_j^2-1)\,u_j^{\,N-i_j}
\prod_{k<i_j}(u_j-w_k)\Bigr]\cdot
\prod_{j>r}u_j^{\,N+1}\cdot
\prod_{\substack{j<k\\ j\le r}}\frac{-u_k}{u_j-u_k}(u_ju_k-1)\cdot
\prod_{r+1\le j<k\le n}\!\!u_ju_k .
\label{eq:B0}
\end{equation}
In particular, if $k\notin\bM$ and $k\ge\max\bI$ then $w_k$ does not occur in
$\At_{\bM,\bI}|_{q=0}$.
\end{lem}
\begin{proof}
Lemma~\ref{lem:q0reg} reduces the claim to $F|_{q=0}$.  In Motegi's Theorem 3.3
the alphabet $\bw$ occurs in exactly three places, and two of them die at $q=0$:
$\prod_{k=i_j+1}^{N}(u_j-qw_k)\to u_j^{N-i_j}$ for $j\le r$, and
$\prod_{k=1}^{N}(u_j-qw_k)\to u_j^{N}$ for $j>r$.  The third,
$\prod_{k=1}^{i_j-1}(u_j-w_k)$ for $j\le r$, carries no $q$ and survives
unchanged.  The remaining factors are $\bw$-free:
$(1-q)(u_j^2-1)\to(u_j^2-1)$,
$\frac{qu_j-u_k}{u_j-u_k}\to\frac{-u_k}{u_j-u_k}$,
$(u_ju_k-q)\to u_ju_k$, and
$\prod_{j>r}(\beta u_j-\gamma)=\prod_{j>r}u_j$ in type $C$, which has been
absorbed into $u_j^{N+1}$ in \eqref{eq:B0}.  This gives
\eqref{eq:q0expand}--\eqref{eq:B0}.  For the last sentence: $w_k$ occurs in
$\mathcal{B}^0$ iff $k<i_j$ for some $j\le r$, i.e.\ iff $k<\max\bI$; and $w_k$
is not among the $\bu=\bw_{\bM}$ because $k\notin\bM$.  Since the condition is
independent of $\sigma$, no $w_k$ can appear through the sum either.
\end{proof}

\begin{cor}[the gap closes on a stratum]\label{cor:stratum}
Suppose $\bM\ge\bI$, $\{a,b\}\cap\bM=\emptyset$ and $\max\bI\le b$.  If
$\At_{\bM,\bI}|_{q=0}\ne0$, then $\ord_{H_{ab}}\At_{\bM,\bI}=0$.
\end{cor}
\begin{proof}
By Lemma~\ref{lem:q0struct}, $w_b$ does not occur in $\At_{\bM,\bI}|_{q=0}$, so
restriction to $H$ --- which is the substitution $w_b\mapsto w_a^{-1}$ --- acts
trivially on it:
$\At_{\bM,\bI}|_{q=0}|_H=\At_{\bM,\bI}|_{q=0}\ne0$.  Apply
Proposition~\ref{prop:q0}.
\end{proof}

\begin{rem}[the residual hypothesis is the wrong way round]\label{rem:q0hyp}
Corollary~\ref{cor:stratum} is conditional on
$\At_{\bM,\bI}|_{q=0}\ne0$, which is \emph{stronger} than the converse of
triangularity ($\bM\ge\bI\Rightarrow\At_{\bM,\bI}\ne0$) that [Q488] already
carries as measured, since $\At|_{q=0}\ne0$ implies $\At\ne0$ but not
conversely.  So the corollary trades one unproved hypothesis for a stronger
one, and is honest progress only to the extent that the $q=0$ object is more
tractable --- which it is, by Lemma~\ref{lem:q0struct}, and which is also where
Motegi identifies $WC|_{q=0}$ with the factorial $GQ$-functions of
Ikeda--Naruse (L555).  I have \emph{not} deep-read Ikeda--Naruse and therefore do
not cite any nonvanishing statement from it; that is the natural next step and
it is a reading task, not a proof task.
\end{rem}

\paragraph{How much is the stratum $\max\bI\le b$?}
Counted combinatorially over the domain of Gap~\ref{gap:nonvanH} (all $a<b$, all
$\bM\ge\bI$ with $\{a,b\}\cap\bM=\emptyset$):
see \S\ref{sec:verif}.  As with [LR]'s minimal stratum, the figure is reported
at several $N$ and the trend, not the first value, is the claim.

\subsection{Two failed attempts, located}
\label{sec:failed}

\paragraph{Attempt 1: a bidegree corner on $\At$.  Refuted by measurement.}
Write $\At_{\bM,\bI}=\sum_{i,j}c_{ij}w_a^iw_b^j$ with
$c_{ij}\in\mathbb{C}(q)[w_k^{\pm1}:k\ne a,b]$; this is legitimate because
$a,b\notin\bM$, so $H$ constrains only $w_a,w_b$.  Then
$\At|_H=\sum_d\bigl(\sum_{i-j=d}c_{ij}\bigr)w_a^d$, and if the maximum of
$i-j$ over the support is attained at a \emph{unique} lattice point then the
corresponding coefficient survives and $\At|_H\ne0$.  The natural candidate for
that point is the corner $(A,B):=(\max_i,\min_j)$.

\textbf{It is empty.}  Measured over the whole domain: $0$ occupied out of $1$
at $N=2$, $0$ out of $12$ at $N=3$, and the same at every case inspected at
$N=4$ (\S\ref{sec:verif}).  The reason is identifiable.  First, I had intended
to read the bidegree off Motegi's symmetrisation formula, where each $w_k$
occurs only in linear factors, one per position $j$, and the positions are not
moved by $\sigma$ --- so each $\sigma$-term has the \emph{same} bidegree,
$\bigl(n-[a\in\bI],\,n-[b\in\bI]\bigr)$.  That is true of $F$ and useless for
$\At$: the individual $\sigma$-terms of $F/\mathcal{E}^C$ are not Laurent
polynomials (the denominator $\prod_{j}(1-qu_jw_a)$ of
$(\mathcal{E}^C)^{-1}$ has $w_a$-degree $n$ and does not divide any single
term), so the bidegree of $\At$ cannot be read term by term.  Second, and
independently fatal, $\At$ carries the explicit factor $(1-qw_aw_b)$ --- a
factor of the $\bw$-product in the denominator of \eqref{eq:EC}, present
because $a,b\in\bM^{\mathrm{c}}$ --- whose two monomials $1$ and $-qw_aw_b$ have
bidegrees $(0,0)$ and $(1,1)$ and hence the \emph{same} value of $i-j$.  Any
such factor ties the maximum, and there is at least one.  A corner argument on
$\At$ cannot work; it would have to be applied to $\At$ divided by its
manifestly $H$-nonvanishing factors, and I did not find a clean statement of
which those are in the unequal-cardinality case.

\paragraph{Attempt 2: the corner at $q=0$.  Reduces to a twisted symmetrisation
sum I cannot evaluate.}
At $q=0$ the first objection vanishes: by Lemma~\ref{lem:q0reg}
$\mathcal{E}^C|_{q=0}=1$, so $\At|_{q=0}$ \emph{is} the symmetrisation sum
\eqref{eq:q0expand} and the bidegree is $\sigma$-uniform.  By
Lemma~\ref{lem:q0struct} the degrees are
$A_0=\#\{i\in\bI:i>a\}$ and $B_0=\#\{i\in\bI:i>b\}$, both minimal degree $0$.
Extracting the corner $(A_0,0)$ --- top in $w_a$, constant in $w_b$ --- the
$w_a$-factors each contribute $-1$ and the $w_b$-factors each contribute $u_j$,
giving
\begin{equation}
[w_a^{A_0}w_b^{0}]\,\At_{\bM,\bI}\big|_{q=0}
=\pm\frac{\prod_ku_k^{-(N+n)+2}}{(n-r)!}
 \sum_{\sigma\in S_n}\sigma\Bigl[
 \Bigl(\prod_{j\le r,\ i_j>a}u_j^{-1}\Bigr)\mathcal{B}^{0,\flat}\Bigr],
\label{eq:twisted}
\end{equation}
where $\mathcal{B}^{0,\flat}$ is \eqref{eq:B0} with the columns $a,b$ deleted
from every product.  When $\bI\cap\{a,b\}=\emptyset$, $\mathcal{B}^{0,\flat}$
\emph{is} the bracket \eqref{eq:B0} for the alphabet $\bw^\flat=(w_k)_{k\ne a,b}$
and the relabelled index set $\bI^\flat$ (the relabelling is order-preserving,
so $\bM\ge\bI$ gives $\bM^\flat\ge\bI^\flat$), and the symmetric monomial
$\prod_ju_j^2$ pulls out of the sum.  \textbf{The obstruction is the remaining
monomial} $\prod_{j\le r,\,i_j>a}u_j^{-1}$: the set $\{j\le r:i_j>a\}$ is the
top $A_0$ \emph{positions}, not a symmetric function of the $u$'s, so it does not
pull out of the $\sigma$-sum, and multiplying a position's block by $u_j^{-1}$
does not shift it to another index's block (it would need one fewer
$(u-w)$ factor, not one lower power of $u$).  So \eqref{eq:twisted} is a
symmetrisation sum twisted by a non-symmetric monomial, and I can neither
identify it with a weight function nor show it nonzero.  When
$\bI\cap\{a,b\}\neq\emptyset$ --- which does occur in the application, e.g.\
$N=3$, $\nu=\{1,2,3\}$, $(a,b)=(1,2)$, $\nu'=\{3\}$, $\lambda=\{1\}$ --- there is
a second obstruction: the block at the position $j$ with $i_j=a$ becomes a
product with \emph{no} hole, which is not an index block of the reduced system
at all.

\paragraph{What I would try next.}
In order: (i) deep-read Ikeda--Naruse on nonvanishing of factorial $GQ$-functions
and see whether it discharges the hypothesis of
Corollary~\ref{cor:stratum} (a reading task --- and the honest statement is that
I currently hold that paper at a level that forbids citing it);
(ii) evaluate \eqref{eq:twisted} in the equal-cardinality case $|\bM|=|\bI|$,
where the $\sigma$-sum has one term and the whole gap is the type $A$ statement
$W_{N,r}(\bw_{\bM}|\bw|\bI)|_H\ne0$ ([LR] \S\texttt{scope} already records
this); (iii) look for the \emph{right} factorisation
$\At=(\text{manifestly }H\text{-nonvanishing})\cdot Q$ in the unequal-cardinality
case, which is what Attempt 1 lacked --- Motegi's divisibility theorems
(Remark~\ref{rem:rootswap}) supply such a factorisation at $|\bM|=|\bI|$, and
their $s>r$ analogue is where I would look.

\section{Computational verification, with controls and plants}
\label{sec:verif}

Pure Python $+$ SymPy over $\mathbb{C}(q)(w_1,\dots,w_N)$ (SageMath is not
installed).  Scripts and logs: \texttt{proofs/code-1017-nonvanH/}.  The engine
\texttt{q484\_engine.py} is carried over; \texttt{eb\_audit.EB} implements
\eqref{eq:EB} and is new.

\paragraph{The normaliser, and that \texttt{EB} is right (\texttt{eb\_audit.py},
\texttt{verify\_all.py}), with two plants.}
Lemma~\ref{lem:EBEC} verified symbolically for every $\bM$ and $r$ at $N=2,3$.
Then \texttt{EB} is checked against Motegi's \emph{independent} theorem
\texttt{typeBmainproperties}(i) (eq.\ \texttt{normalizationWB} L1386), i.e.\
against the type $B$ diagonal root-product formula, which is a different theorem
from the normalisation factor \eqref{eq:EB}:
\begin{center}
\begin{tabular}{@{}lrr@{}}
\toprule
 & $N=2$ & $N=3$ \\
\midrule
\texttt{EB} as implemented: match / cases & $4/4$ & $8/8$ \\
\textbf{P-A1} inclusive $\bw$-diagonal (the type $C$ convention): match & $0/4$ & $0/8$ \\
\textbf{P-A2} drop $\prod_j(1-qu_j)$: match & $1/4$ & --- \\
\bottomrule
\end{tabular}
\end{center}
P-A1 is the sharp plant: it is exactly the error the engine was making, and it
fails on every case.  P-A2 is \emph{informatively spared} on one case, namely
$\bI=\emptyset$ with $n=0$, where its product is empty and the plant is the
identity by construction.  The two plants sit in different blocks of
\eqref{eq:EB} --- the $\bw$-diagonal and the linear $\bu$-product --- which is
the point of naming them separately.

\paragraph{The multiplier (\texttt{eb\_rho.py}, \texttt{rhoB\_N4.py}), two arms.}
Over all $(a<b,\ \bM\supseteq\{a,b\},\ \bI)$, restricted to $H$, reporting only
the \emph{non-degenerate} domain (some side $\ne0$ on $H$; a pair with both
sides zero cannot distinguish any multiplier):
\begin{center}
\begin{tabular}{@{}lrrr@{}}
\toprule
object tested & domain & $\rho^C$ holds & $\rho^B$ holds \\
\midrule
$\widetilde{WC}=WC/\mathcal{E}^C$, $N=2$ & $1$ & $1/1$ & $0/1$ \\
$\widetilde{WC}=WC/\mathcal{E}^C$, $N=3$ & $12$ & $12/12$ & $0/12$ \\
$\widetilde{WB}=WB/\mathcal{E}^B$, $N=2$ & $1$ & $0/1$ & $1/1$ \\
$\widetilde{WB}=WB/\mathcal{E}^B$, $N=3$ & $12$ & $\mathbf{0/12}$ & $\mathbf{12/12}$ \\
\bottomrule
\end{tabular}
\end{center}
Each candidate holds on its own type and fails on the other, and the two
candidates disagree on the \emph{whole} domain, so the test is discriminating
rather than merely green: it is evaluated on the disagreement set, and the
disagreement set is not empty.

\paragraph{Reproducing the retracted number (\texttt{rhoB\_N4.py}), closing the
derivation loop.}
The proof of Theorem~\ref{thm:rhoB} routes through $\Lambda=WB/\mathcal{E}^C$,
so it needs [LR]'s measurement as an input rather than as a corroboration.  That
measurement reproduces: $\Lambda$ collapses with $\rho^C$ in $1/1$
non-degenerate of $4$ pairs at $N=2$ and $12/12$ non-degenerate of $48$ pairs at
$N=3$ --- the same $48$ and $12$ as [LR] reports, so the retracted number was a
correct measurement of the wrong object.  Plants in the multiplier block:
$\rho\mapsto\rho(1+q)$ and $\rho\mapsto1$ both fail $0/12$, as required.  Hence
$\rho^B$ is reached twice independently: by \eqref{eq:Gratio} from the
reproduced $\Lambda$-measurement plus Lemma~\ref{lem:EBEC}, and by direct
measurement of $\widetilde{WB}$ in the table above.  The $N=4$ arm for
$\widetilde{WB}$ had not returned at write-up and is marked
\emph{pending}: a claim arm without its plants, or a table without its widest
domain, is reported as pending rather than assumed.

\paragraph{The dictionary (symbolic).}
$\rho^X-1-q\,w_a^{h}([K^X(w_a)]_0^{\,1})^2/\bigl((1-qw_a^{h})(1-qw_a^{-h})\bigr)$
simplifies to $0$ for $X=C$ ($h=2$) and $X=B$ ($h=1$); and
$w_a^{h/2}[K^X(w_a)]_0^{\,1}-(w_a^{h/2}-w_a^{-h/2})=0$ in both.  The values in
the disagreement table of \S\ref{sec:dict} were computed, not asserted.

\paragraph{Attempt 1, refuted (\texttt{corner.py}).}
Over the whole domain of Gap~\ref{gap:nonvanH}: the bidegree corner
$(\max_i,\min_j)$ of $\At_{\bM,\bI}$ in $(w_a,w_b)$ is \textbf{occupied in $0$
cases and empty in $1$ of $1$ at $N=2$, $12$ of $12$ at $N=3$}, and likewise at
every $N=4$ case inspected.  This is a measurement that refuted a route before
any time was spent proving it, and it is recorded as the only outcome of
Attempt 1.

\paragraph{Lemma~\ref{lem:q0struct} (\texttt{q0\_struct.py},
\texttt{q0b.py}), with a plant, and a pre-registered claim that was wrong.}
First, $\mathcal{E}^C|_{q=0}=1$ for every $\bM$ and $r$ at $N=2,3,4$.  Then the
occurrence question.  I pre-registered the \emph{iff}: ``$w_k$ occurs in
$\At_{\bM,\bI}|_{q=0}$ if and only if $k<\max\bI$''.  \textbf{The iff is false}:
it fails at $N=3$, $\bM=\{2,3\}$, $\bI=\{1,2\}$, $k=1$, where $w_1$ is absent
although $1<\max\bI=2$, and again twice at $N=4$.  The direction
Lemma~\ref{lem:q0struct} actually asserts and
Corollary~\ref{cor:stratum} actually uses --- $k\ge\max\bI\Rightarrow w_k$
absent --- holds in $6/6$ at $N=2$, $27/27$ at $N=3$ and $110/110$ at $N=4$.
So the lemma stands as
stated and there is additional vanishing beyond it that I did not predict.
The \textbf{plant} was the same claim with $k\le\max\bI$; it disagrees with the
measurement $1$ time at $N=2$ and $7$ times at $N=3$, which is what shows the
two thresholds are distinguishable by this scan at all ($8$ disagreements at
$N=4$).  Recording this
separately matters: a pre-registered \emph{iff} that fails in its harmless
direction is a bookkeeping error in the prediction, not in the lemma, and
collapsing the two directions would have reported either a false lemma or a
false green.

\paragraph{The stratum of Corollary~\ref{cor:stratum} (\texttt{q0b.py}).}
Counted combinatorially over all $(a<b,\ \bM\ge\bI,\ \{a,b\}\cap\bM=\emptyset)$:
\begin{center}
\begin{tabular}{@{}lrrrrrrr@{}}
\toprule
$N$ & $2$ & $3$ & $4$ & $5$ & $6$ & $7$ & $8$ \\
\midrule
domain              & $1$ & $12$ & $94$ & $610$ & $3558$ & $19376$ & $100588$ \\
$\max\bI\le b$      & $1$ & $11$ & $77$ & $444$ & $2306$ & $11241$ & $52573$ \\
\%                  & $100.0$ & $91.7$ & $81.9$ & $72.8$ & $64.8$ & $58.0$ & $52.3$ \\
\bottomrule
\end{tabular}
\end{center}
\textbf{Decreasing}, and reported over seven consecutive $N$ rather than from
its head: [LR]'s minimal stratum went
$100\%\to92\%\to78\%\to61\%\to44\%$ and its first two values read as a closure.
Corollary~\ref{cor:stratum} is a partial closure on a stratum that is
plausibly asymptotically small --- the figures are consistent with a limit at or
below $\tfrac12$ and I make no claim about the limit --- so it does not reduce
Gap~\ref{gap:nonvanH} to a boundary case.

\section{Scope}
\label{sec:scope}

\paragraph{Retracted.}
[LR] Theorem \texttt{collapse}, final sentence (``the same $\rho$ holds in type
$B$'') and its $48/48$; [Q484] Theorem A, type $B$ arm ($48/48$).  Both measured
$WB/\mathcal{E}^C$.  The type $B$ \emph{statement} of
Theorem~\ref{thm:rhoB} replaces the former; the latter awaits re-measurement.

\paragraph{Unaffected.}
Everything in [LR] and [Q484] stated for type $C$.  $\mathcal{E}^C$ is the
correct normaliser for type $C$ and the engine's type $C$ arm is checked against
Motegi's independent diagonal root-product formula.  In particular
[LR] Theorems \texttt{easy}, \texttt{simple}, \texttt{recursion},
\texttt{created} and Proposition \texttt{rho} stand as type $C$ statements ---
but [LR]'s remark that they ``hold in type $B$ to whatever extent those facts
do'' must now be read with Theorem~\ref{thm:rhoB}: the facts (T), (R), (D), (S),
(A) have type $B$ analogues with the short-root substitution, and the type $B$
versions of those theorems carry $\rho^B$, not $\rho^C$.  I have not worked them
out.

\paragraph{Open.}
Gap~\ref{gap:nonvanH} (\texttt{nonvanH}), now with a reduction to $q=0$
(Proposition~\ref{prop:q0}), an explicit $q=0$ structure
(Lemma~\ref{lem:q0struct}), a conditional closure on $\max\bI\le b$
(Corollary~\ref{cor:stratum}), and two routes closed off with reasons
(\S\ref{sec:failed}).  [LR] Gap \texttt{prop} (\texttt{propagation}), untouched,
and still the bulk of the problem.  Coordinate 3 of the boundary dictionary
(Remark~\ref{rem:coord3}): the mechanism.

\paragraph{Q500.}
The brief gated the boundary reading on a slot check with three pre-registered
outcomes.  The verdict is \textbf{0A, same slot} --- the $K$-matrices of
$\At_{\bM,\bI}$ are indexed by the specialisation set $\bM$, a collapse deletes
exactly two of them, and the charge-violating entry is a divisor of the weight
function in Motegi's own divisibility theorems --- reached by refuting the
brief's own version of the claim (wrong paper, wrong entry, and a $\rho$ that
turned out to be type-dependent).  It is promoted from shape to mathematics with
a named mechanism gap, not to a theorem.

\section*{Source hygiene}

Every fact attributed to Motegi is from arXiv:2610.15344, read first-hand in the
e-print source \texttt{typebcweightfunctions\_arxivversionfinal-3.tex}, with
line numbers given at the point of use:
\eqref{eq:Ktri} L332--345; \eqref{eq:Kentries} and \eqref{eq:KCKB} L352--375;
\eqref{eq:WCWB} L489--505; \eqref{eq:EC} L1051--1062; \eqref{eq:EB}
L1361--1367; the two divisibility theorems L1086--1089 and L1394--1402; the
type $B$ diagonal formula L1386; the symmetrisation formula (Motegi Thm 3.3)
L491--504; the $q=0$ identification with factorial $GQ$/$GB$-functions L555.
The $K$-matrix of arXiv:1703.07924 is quoted from
\texttt{arxiv-1703.07924-source.TEX} L300--321, and is quoted only in order to
record that it is \emph{not} the $K$-matrix of the paper under discussion.
Ikeda--Naruse is named in Remark~\ref{rem:q0hyp} as a reading task and no claim
rests on it.  No claim above rests on a document whose author is also its
subject.

\end{document}
